% ============================================================================
% PART 01 - INTRODUCTION AND MOTIVATION
% ============================================================================

\labelsec{intro_motivation}

\section{The DNA Geometry Problem}

DNA structure determines function. This principle, foundational to molecular biology, has driven decades of structural investigation from fiber diffraction to modern cryo-EM. Yet our existing descriptive frameworks remain incomplete.

\claim{Current description methods capture helical orientation (twist, roll, tilt) but not cross-sectional geometry.}

The canonical helical parameters \citep{Dickerson1989}, derived from classical work \citep{Watson1953}, describe the relative positions and orientations of successive base pairs. Tools like Curves+ \citep{Lavery1989} and 3DNA \citep{Lu2003} compute these parameters at high precision. These parameters answer: \textit{How is this base pair rotated and displaced relative to the previous one?}

But they do not answer: \textit{What is the \textbf{shape} of the cross-section perpendicular to the helix axis? How does local bending deform the DNA? Where is the backbone most strained?}

\section{Prior Approaches and Limitations}

\subsection{Curves+ Axis Curvature}

Curves+ \citep{Lavery1989} computes helical axis curvature $\kappa = 1/R$, where $R$ is the instantaneous radius of curvature. This describes the 1D path of the helix axis through space.

\textbf{Limitation:} Global property (axis bending) does not characterize local distortion. A structure could have moderate axis curvature but severe cross-sectional deformation.

\subsection{3DNA and Helical Parameters}

3DNA \citep{Lu2003} provides refined helical parameters including rise, twist, roll, tilt, and inclination.

\textbf{Limitation:} These describe relative base-pair step geometry, not absolute cross-sectional shape. They cannot distinguish between bending caused by asymmetric base stacking vs. bending caused by protein-induced asymmetry.

\subsection{Sequence-Based Prediction (Canal, NuPoP)}

Tools like Canal \citep{Vlahovicek2003} and NuPoP \citep{Xi2010} predict DNA geometry from sequence using trained models.

\textbf{Limitation:} Predictions rather than measurements. Cannot validate predictions without calculating actual geometry from structures.

\subsection{Rigidity and Flexibility Parameters}

Classical approach characterizes DNA through persistence length and bending rigidity.

\textbf{Limitation:} Bulk properties; cannot measure local shape at single-basepair resolution.

\section{The Missing Gap}

\begin{highlight}
    No existing tool systematically measures \textbf{cross-sectional geometry} of DNA at basepair resolution across diverse structural contexts.
\end{highlight}

This gap has profound implications:

\begin{enumerate}
    \item \textbf{Nucleosome wrapping:} We know DNA wraps 1.65 turns around histones, but cannot quantify how the cross-section deforms during wrapping.
    
    \item \textbf{Protein-DNA recognition:} Transcription factors recognize DNA partly through shape complementarity \citep{Rohs2010}, yet we lack quantitative shape descriptors.
    
    \item \textbf{Sequence-structure relationships:} Why do A-tracts resist bending? Not fully explained by helical parameters alone.
    
    \item \textbf{Functional prediction:} Current DNA classification relies on sequence features or global properties, missing local geometric signatures.
\end{enumerate}

\section{The Shape-Height Transform Solution}

We propose the Shape-Height Transform (SHT), a rigorous geometric framework for measuring DNA cross-sectional properties at basepair resolution.

\subsection{Core Concept}

Given a PDB structure:

\begin{enumerate}
    \item Extract DNA backbone coordinates
    \item Compute the principal axis through the structure
    \item Define height $h$ as position along this axis
    \item At each height, establish a local coordinate frame perpendicular to the axis
    \item Compute moments of atomic distribution in cross-sectional plane
    \item Derive three orthogonal geometric units from these moments
\end{enumerate}

\subsection{Three Orthogonal Units}

\begin{keyresult}
    \textbf{Vij} (local curvature, $v_0$): Measures how strongly the helix axis curves. Quantifies DNA bending at each basepair.
    
    \textbf{Shambhian} (cross-sectional flatness, $\eta_r$): Measures deviation from circular cross-section. Characterizes asymmetry in helix geometry.
    
    \textbf{Mamton} (surface roughness, $\Omega$): Measures skewness of atomic distribution. Captures asymmetric distortion and wrapping-induced deformation.
\end{keyresult}

\subsection{Why These Three Units?}

These units are:

\begin{itemize}
    \item \textbf{Mathematically orthogonal:} Derived from different moments (0th, 2nd, 3rd) of the distribution
    \item \textbf{Physically meaningful:} Each corresponds to a real geometric property
    \item \textbf{Computationally tractable:} Can be calculated from all PDB structures
    \item \textbf{Dimensionally clear:} Have well-defined units and ranges
    \item \textbf{Functionally discriminative:} Each captures different aspects of DNA geometry
\end{itemize}

\section{Research Questions}

This monograph addresses three primary research questions:

\begin{enumerate}
    \item \textbf{Q1: Can we measure DNA cross-sectional geometry reliably?}
    
    We develop the SHT mathematical framework and validate each unit against known structures and established tools.
    
    \item \textbf{Q2: Are these geometric measurements independent or coupled?}
    
    We test pairwise correlations between units and investigate mechanistic coupling.
    
    \item \textbf{Q3: Can geometry predict DNA function?}
    
    We build a predictive classifier and validate on independent structures, including blind predictions.
\end{enumerate}

\section{Document Roadmap}

\textbf{Part I: Theory} (Chapters 1-3)
\begin{itemize}
    \item Chapter 1: Introduction (this chapter)
    \item Chapter 2: Mathematical framework (height function, tangent frames, moments)
    \item Chapter 3: Geometric unit derivations (formulas for Vij, Shambhian, Mamton)
\end{itemize}

\textbf{Part II: Methods} (Chapters 4-5)
\begin{itemize}
    \item Chapter 4: Algorithm and software implementation
    \item Chapter 5: Quality control and edge cases
\end{itemize}

\textbf{Part III: Validation} (Chapters 6-7)
\begin{itemize}
    \item Chapter 6: Unit-by-unit validation with test cases
    \item Chapter 7: Cross-unit relationships and coupling
\end{itemize}

\textbf{Part IV: Results} (Chapters 8-9)
\begin{itemize}
    \item Chapter 8: Functional classification and prediction
    \item Chapter 9: Biological insights and applications
\end{itemize}

\textbf{Part V: Conclusions} (Chapters 10-11)
\begin{itemize}
    \item Chapter 10: Summary of evidence
    \item Chapter 11: Limitations and future work
\end{itemize}

\section{Key Findings Preview}

This monograph establishes:

\begin{keyresult}
    \begin{enumerate}
        \item Each geometric unit is mathematically valid and measures its target property
        \item Vij and Mamton exhibit strong coupling ($r = 0.83$, $p < 10^{-16}$)
        \item Four distinct geometric states exist in the structural database
        \item Geometric signatures predict functional class with 91\% accuracy
        \item Blind prediction on 1F66 (nucleosome) validated via PDB metadata
    \end{enumerate}
\end{keyresult}

All findings discussed here are substantiated by quantitative evidence ($p < 0.01$, large effect sizes) presented in detail in subsequent chapters.
